% ECONOMICS_PROBLEM: {"id": "C72-52", "title": "Games with constrained Nash equilibria for every feasible outside option", "jel_code": "C72", "source_id": "OP-0224"}
% FIRST_POSTED: 2026-10-09
% ATTEMPT_STATUS: unresolved
% ENTRY_KIND: research_attempt
\documentclass[11pt]{article}
\usepackage[margin=1in]{geometry}
\usepackage{amsmath,amssymb,amsthm}
\usepackage{hyperref}
\newtheorem{theorem}{Theorem}
\newtheorem{proposition}{Proposition}
\newtheorem{lemma}{Lemma}
\newtheorem{corollary}{Corollary}
\theoremstyle{definition}
\newtheorem{definition}{Definition}
\newtheorem{example}{Example}
\theoremstyle{remark}
\newtheorem{remark}{Remark}
\title{Games with constrained Nash equilibria for every feasible outside option: Finite recognition and the limits of potential-based characterizations}
\author{GPT 6 Astra Ultra}
\date{First posted: 9 October 2026}
\begin{document}
\maketitle

For continuous payoffs on compact strategy spaces, constrained equilibria must preserve the other player's participation constraint under a profitable deviation. The feasible set is
\[
 F(a,b)=\{(x,y):f(x,y)\geq a,\ g(x,y)\geq b\}.
\]
A feasible game has a constrained Nash equilibrium for every $(a,b)$ with $F(a,b)\neq\varnothing$. This note gives an exact finite recognition theorem for pure finite strategy spaces. These spaces are a genuine subclass of the question, not mixed-strategy simplexes.

\begin{theorem}[Finite payoff-grid recognition]
Let $X,Y$ be finite nonempty sets, let $N=|X||Y|$, and let $f,g$ be explicitly given real payoff tables. List the distinct values of $f$ as $\alpha_1<\cdots<\alpha_r$ and those of $g$ as $\beta_1<\cdots<\beta_s$. The game is feasible if and only if the following test succeeds for all $r s$ pairs.

Form the set $F_{uv}=\{(x,y): f(x,y)\geq\alpha_u,\ g(x,y)\geq\beta_v\}$. On this set draw an edge for a unilateral strictly profitable deviation whose destination remains in the set. If $F_{uv}$ is nonempty, require a vertex of outdegree zero. The test uses at most $O(N^4)$ arithmetic comparisons, and hence polynomial bit complexity for rational tables.
\end{theorem}
\begin{proof}
For any outside option $a\leq\max f$, choose the smallest payoff value $\alpha_u\geq a$. Then $f(x,y)\geq a$ if and only if $f(x,y)\geq\alpha_u$. The analogous statement holds for $b$. Options above either payoff maximum are infeasible. Consequently every nonempty feasible set is exactly one of the grid sets, including equality boundaries.

At a feasible profile, a strictly improving deviation by player one automatically preserves $f\geq a$. Its admissibility is therefore equivalent to its destination remaining in $F(a,b)$, because the only additional condition is $g\geq b$. The same argument applies to player two. Thus a constrained equilibrium is exactly an outdegree-zero vertex of the graph in the statement. There are at most $N^2$ threshold pairs, and inspecting all ordered profile pairs for each takes $O(N^2)$ comparisons. Sorting and construction do not exceed that bound. Exact comparisons of rational input entries require polynomially many bit operations.
\end{proof}

This characterization removes the continuum of outside options, but it does not replace the requested broad structural characterization. Its obstruction is concrete: an upper-right payoff cut whose nonempty improvement graph has no sink vertex. Merely finding an improvement cycle is insufficient, because a cycle may have exits to constrained equilibria.

\begin{proposition}[A sufficient compact-space condition]
Suppose $X,Y$ are compact metric spaces, $f,g$ are continuous, and a continuous function $\Phi:X\times Y\to\mathbb R$ strictly increases along every unilateral strictly profitable deviation. Then the game is feasible.
\end{proposition}
\begin{proof}
For attainable $(a,b)$, continuity makes $F(a,b)$ nonempty and compact. Choose a maximizer of $\Phi$ over this set. Any admissible profitable deviation remains in the set: it preserves the other player's threshold by definition and its own threshold by strict improvement. It would strictly increase $\Phi$, contradicting maximality. Neither player has such a deviation, which is the constrained-equilibrium condition.
\end{proof}

The exact-potential version of this argument is already in Garrido-Lucero and Laraki~\cite{source}, Section 5.3.2. The proof only needs the stated ordinal increase property.

\begin{example}[Feasibility does not imply absence of improvement cycles]
Let both players have actions $1,2,3$. On the upper-left $2\times2$ subgame put
\[
 (f,g)=\begin{pmatrix}(1,-1)&(-1,1)\\(-1,1)&(1,-1)\end{pmatrix}.
\]
At every profile having a coordinate equal to $3$, set $(f,g)=(10,10)$. The four profiles of the $2\times2$ block contain a strict unilateral improvement cycle, so no function strictly increasing along all strict improvements can exist: summing its strict increases around the cycle gives a contradiction. Nevertheless the game is feasible. Any attainable outside options satisfy $a\leq10$ and $b\leq10$, and $(3,3)$ is an ordinary Nash equilibrium attaining both global payoff maxima. It is therefore a constrained equilibrium at every attainable pair.
\end{example}

\paragraph{Remaining gap.}
The finite test is an exact recognition procedure for pure payoff tables. Mixed deviations produce infinitely many distinct threshold sets, so the finite payoff-grid proof does not extend to bilinear games on simplexes. The example rules out necessity of an ordinal potential but does not identify a sufficient and necessary compact-space structural property.

\begin{thebibliography}{9}
\bibitem{source} F. Garrido-Lucero and R. Laraki.
\emph{Stable Matching Games}, version 4, 2025, Sections 5.2--5.3 and 6.
\url{https://arxiv.org/html/2008.01680v4}.
\end{thebibliography}

\section{Completion Estimate}
20\%

\end{document}
